\documentclass[10pt,a4paper]{amsart}
\usepackage[margin=19mm]{geometry}
\usepackage{amsmath,amssymb,mathtools}
\usepackage[scr=rsfs,cal=euler]{mathalfa}
\usepackage{xfrac}
\usepackage[unicode,hidelinks,bookmarks=false]{hyperref}
\newcommand{\ie}{i.e.\ }
\newcommand{\cf}{cf.\ }
\newcommand{\resp}{respectively\ }
\newcommand{\Subsection}{\subsection}
\providecommand{\nobreakdash}{\nobreak\hbox{-}}
\setlength{\emergencystretch}{2em}
\multlinegap0pt
\usepackage{amssymb}
\usepackage{mathtools}
\usepackage{bm}
\usepackage{enumerate}
\usepackage{bbm}
\newtheorem{proposition}{Proposition}
\newcommand \PPb{\mathbb{P}}
\newcommand \Val{{\rm{CoV}}}
\newcommand \rng{{\rm{ rng}}}
\title{A logic of co-valuations}
\author{Maciej Malicki}
\date{Source: arXiv:2511.01011v2 (CC BY 4.0)}
\begin{document}
\maketitle

\noindent\emph{Source and licence.} A logic of co-valuations. Version arXiv:2511.01011v2 (CC BY 4.0), distributed by arXiv under the Creative Commons Attribution 4.0 International licence, http://creativecommons.org/licenses/by/4.0/, as recorded in the arXiv licence metadata for this exact version and shown on the abstract page https://arxiv.org/abs/2511.01011v2. Full licence notice: \texttt{licenses/arxiv-2511.01011-CC-BY-4.0.txt}.\par\smallskip

\noindent\textit{Editorial scope.} Counted unit: the article's proposition pr:FormulaProp (source0.tex lines 768--779). The article's complete original proof of it is delivered with it (source0.tex lines 781--805, one \texttt{proof} environment of 25 lines). No mathematics is written, repaired or re-derived: the statement and the proof are the article's own lines, in the article's own notation, and no retained line is substituted or re-flowed.  No further source-local context is needed: every cross-reference of every retained line resolves inside the two delivered spans.  Nothing was omitted from any retained span, so the package carries no omission record.  No retained line contains a citation, so the wrapper carries no bibliography and no bibliography entry.  No retained line contains a figure environment, an image-inclusion command or drawing code, so no image is delivered and the package declares no asset dependency.  Editorial formatting only, in addition: an AMS \texttt{amsart} preamble that restates the article's own declarations and macros used by the retained lines, namely the environments \texttt{proposition} on separate counters, together with the standard packages those lines need.  The retained environments print in this document as Proposition 1 in order of appearance, whereas the article prints the same items with the numbering of its own full document.  The complete native source of the article as distributed in the arXiv e-print for this exact version is bundled unchanged as \texttt{sources/188290/source0.tex}.
\begin{proposition}
	\label{pr:FormulaProp}	
	Let $(M,V)$ be a compact structure, let $X$ be a context, let $x_0 \in X$, $X_0 \subseteq X$, where $|X_0|>1$,  and $v \in V^X$. The following properties of $v$ can be expressed by a single formula:
	
	\begin{enumerate}
		\item $\bigcap [v]_{X_0}= \emptyset$, 
		\item $[v]$ is a partition,
		\item $[x_0]_v$ is a singleton.
		%\item $\PPb(v)$ is a chain, i.e., there is an enumeration $x_0, \ldots, x_n$ of $X$ such that $[x_i]v \cap [x_j]v \neq \emptyset$ iff $|i-j|=1$. 
		%\item $w$ is a star-refinement of $v$.
	\end{enumerate}
\end{proposition}		

\begin{proof}
	
	We show (1). Define $Y=X \cup \{y_0\}$, where $y_0 \not \in X$, and $\geq \in \Val(X,Y)$ by putting $x \geq y$ iff $x=y$ or ($x \in X_0$ and $y=y_0$). Define
	%
	\[ \text{HasEmptyIntersection}_{X_0}(X):= \forall_{\geq} \bot.\]
	%
	Clearly, if $[v]_{X_0}$ has empty intersection, then $$M \models_v \text{HasEmptyIntersection}_{X_0}(X).$$ To see the other direction, suppose that $c=\bigcap [v]_{X_0 } \neq \emptyset$. Fix a minimal open cover $C$ of $M$ extending $\{c\}$,  and  let $w' \in V^Z$ be a fragmentation of $v$ such that $[w']$ refines $C$. Fix $y_0 \in Z$ such that $[y_0]_{w'} \subseteq c$. Note that, by minimality of $[v] $, for every $x \in X$, there is $z \in Z$ such that $z \neq y_0$, $[x]_v \supseteq [z]_{w'}$, and $[x']_v \supseteq [z]_{w'}$ implies $x=x'$.  Finally, let $\succeq \in \Val(Y,Z)$ be a pattern defined by $y \succeq z$ iff $y=z=y_0$ or ($z \neq y_0$ and  $[x]_v \supseteq [z]_{w'}$). Then $w=C_\succeq (w')$ witnesses that $M  \models_v   \exists_{\geq} \top$.
	
	
	%and fix the largest $X_ 1 \supseteq X_0$ such that $p=\bigcap [v]_{X_1}=\bigcap [v]_{X_0}$. Select a finite $M_1 \subseteq M$ such that, for $x \in X$, $M_1 \subseteq [x]$ iff $x \in X_1$. Select a finite cover $C$ of $M$ such that $p \in C$, and $q \cap M_1  \neq \emptyset$ iff $q=p$, for $q \in C$.  Let $w' \in V$ be a fragmentation of $v$ such that $[w']$ refines $C$. Finally, let $w$ be a consolidation of $w'$ deteremined by $Y_{X_1}$. Then $w$ witnesses that $M  \models_v   \exists_{\geq Y_{Z}} \top$.
	
	In order to show (2), we define 
	\[\text{IsAPartition}(X):= \bigwedge_{\{x_1,x_2 \in X: x_1 \neq x_2\}} \text{HasEmptyIntersection}_{\{x_1,x_2\}}(X).\]
	%
	Finally, we show (3). Let $Y=X \cup \{y_0,y_1\}$, where $y_0,y_1\ \not \in X$, $y_0 \neq y_1$.  Define $\geq \in \Val(X,Y)$ by putting $x \geq y$ iff $x=y$ or ($x=x_0$ and  ($y=y_0$ or $y=y_1$)). Let
	% 
	\[\text{IsASingleton}_{x_0}(X):=\forall_{\geq} \bot .\]
	
	Clearly, if $[x_0]_v$ is a singleton, then $$M \models_v \text{IsASingleton}_{x_0}(X).$$ To see the other direction, fix distinct $a_0, a_1 \in [x_0]_v$. Fix open $c_0, c_1 \subseteq [x_0]_v$ such that $a_0 \in c_0 \setminus c_1$, $a_1 \in c_1 \setminus c_0$,  a minimal open cover $C$ of $M$ extending $\{c_0, c_1\}$ such that $a_0, a_1 \not \in \bigcup C \setminus \{c_0,c_1\}$,  and  let $w' \in V^Z$ be a fragmentation of $v$ such that $[w']$ refines $C$. Fix $y_0, y_1 \in Y$  such that $a_0 \in [y_0]_{w'}$, $a_1 \in [y_1]_{w'}$; in particular $y_0 \neq y_1$.  Finally, let $w$ be a consolidation of $w'$ determined by $Y$ as in (1). Then $w$ witnesses that $M  \models_v   \exists_{\geq} \top$.
	%
	%	\item  Let $n$ be the size of $X$.
	%
	%	\[ IsAChain(X)=\bigvee_{ \bar{y} \in X^n: \rng(\bar{y})=X } \bigwedge_{i,j \leq n : |i-j|>1} \text{HasEmptyIntersection}_{y_i,y_j}(X) \wedge \]
	%	\[ \wedge  \bigwedge_{i,j \leq n : |i-j|=1} \neg \text{HasEmptyIntersection}_{y_i,y_j}(X).  \]
\end{proof}

\end{document}
